% -------------------------------------------------
% Capítulo: Diseño
% Orden del director (19-sep-2026): diseño arquitectónico; qué herramientas se
% decidieron y por qué estas y no otras; diseño de la interfaz de usuario.
% Los patrones de diseño del código van en el capítulo de Implementación.
% Versiones y nombres verificados contra composer.json, package.json,
% migraciones y controladores (25-sep-2026).
% -------------------------------------------------

Este capítulo explica cómo se diseñó Pymetory a partir de los requerimientos del Capítulo \ref{chap:requerimientos}: la arquitectura general del sistema, las herramientas elegidas y las razones de cada elección, el modelo de datos, el diseño del módulo de consulta en lenguaje natural y el diseño de la interfaz de usuario.

\section{De los requerimientos a los componentes}
\label{sec:dis_trazabilidad}

Cada requerimiento funcional implementado se asignó a uno o varios componentes del sistema. La Tabla \ref{tab:dis_trazabilidad} muestra esa correspondencia; los nombres de clases y comandos son los del código fuente.

{\small
\begin{longtable}{|c|p{5.2cm}|p{7cm}|}
\caption{Correspondencia entre requerimientos funcionales y componentes.}
\label{tab:dis_trazabilidad}\\
\hline
\textbf{RF} & \textbf{Requerimiento (resumen)} & \textbf{Componentes} \\
\hline
\endfirsthead
\hline
\textbf{RF} & \textbf{Requerimiento (resumen)} & \textbf{Componentes} \\
\hline
\endhead
RF-01, RF-02 & Catálogo de materiales e ingreso de lotes & \texttt{InventoryController}, modelos \texttt{Material} y \texttt{Lote} \\ \hline
RF-04 & Despacho FEFO & \texttt{ConsumptionController}, \textit{scope} \texttt{Lote::despachables} \\ \hline
RF-05 & Kardex & Modelo \texttt{Movimiento}, tabla \texttt{movimientos} \\ \hline
RF-06, RF-18 & Roles y usuarios & \textit{Middleware} \texttt{CheckRole}, \texttt{UserController}, \texttt{AuthController} \\ \hline
RF-07, RF-08 & Alertas de stock bajo y de vencimiento & Comandos \texttt{alerts:low-stock} y \texttt{alerts:fefo}, programados cada hora \\ \hline
RF-09, RF-10 & Punto de reorden y consumo por periodo & Servicio \texttt{ProyeccionInventario}, \texttt{ProyeccionController}, vista Reabastecimiento, intenciones \texttt{forecast} y \texttt{consumption} del asistente \\ \hline
RF-17 & Reportes & \texttt{ReportController}, \texttt{InventoryController} (exportación PDF, CSV y Excel) \\ \hline
RF-11 & Conciliación & \texttt{InventoryController::adjust} \\ \hline
RF-12 & Consulta en lenguaje natural & \texttt{ChatLLMController} \\ \hline
RF-13 & Etiquetas y escaneo & \texttt{LabelController}, \texttt{QRScanController} \\ \hline
RF-14 & Recepción de órdenes de compra & \texttt{PurchaseOrderController} \\ \hline
RF-16 & Bodegas y transferencias & \texttt{InventoryController}, \texttt{TransferController} \\ \hline
\end{longtable}
}

\section{Arquitectura del sistema}
\label{sec:dis_arquitectura}

Pymetory es una aplicación web monolítica organizada según el patrón Modelo-Vista-Controlador (MVC) que propone el framework Laravel \cite{laraveldocs}. La capa de presentación se construye con React y se conecta con los controladores de Laravel mediante Inertia.js \cite{inertiajs}: el controlador responde con el nombre de una página y los datos que necesita, e Inertia se encarga de renderizar el componente React correspondiente en el navegador. De esta forma la aplicación se comporta como una aplicación de una sola página (SPA) sin que sea necesario diseñar y mantener una API REST separada para las vistas principales. Algunas funciones puntuales (usuarios, etiquetas, órdenes de compra, proveedores de modelos de lenguaje) sí exponen rutas JSON bajo el prefijo \texttt{/api}, que consumen los componentes React con Axios.

La Figura \ref{fig:dis_arquitectura} resume las capas y el recorrido de una solicitud.

\begin{figure}[H]
\centering
\footnotesize
\begin{verbatim}
 Navegador (React 19 + Inertia.js)
        |  solicitud HTTP (sesión)
        v
 Laravel 13 - middleware: auth, verified, role:admin (CheckRole)
        |
        +--> Controladores de inventario --> Modelos Eloquent --> MySQL 8
        |      (Inventory, Consumption,        (Material, Lote,
        |       Transfer, Report, Label...)     Movimiento, Bodega...)
        |
        +--> ChatLLMController (RAG)
               1. intencion()        intención (modelo; patrones de reserva)
               2. buildRagContext()  lectura de MySQL
               3. llamada al proveedor de modelo de lenguaje
                  (Ollama | OpenCode | API compatible OpenAI; respaldos)
               4. respuesta + registro en chat_histories
\end{verbatim}
\caption{Arquitectura en capas de Pymetory y flujo de una solicitud.}
\label{fig:dis_arquitectura}
\end{figure}

Dos decisiones de diseño se derivan directamente de las reglas de la dirección del trabajo de grado:

\begin{itemize}
    \item \textbf{El módulo de lenguaje natural es de solo lectura sobre el inventario.} El \texttt{ChatLLMController} consulta las tablas del inventario para construir el contexto, pero no crea, modifica ni elimina lotes o movimientos; solo escribe en su propio historial de conversación. Así, un error del modelo de lenguaje puede producir una respuesta equivocada, pero no puede alterar las existencias.
    \item \textbf{El control de acceso se decide en el servidor.} Las rutas administrativas se protegen con el \textit{middleware} \texttt{role:admin}; la interfaz oculta opciones según el rol, pero la autorización real ocurre en Laravel, antes de llegar al controlador.
\end{itemize}

\section{Selección de herramientas}
\label{sec:dis_herramientas}

La selección de herramientas respondió a cuatro criterios propios del proyecto: (1) costo de licencia nulo, porque el sistema está pensado para PYMEs que no pueden pagar suscripciones; (2) un solo desarrollador, lo que favorece herramientas que resuelven de fábrica autenticación, migraciones y pruebas; (3) experiencia previa del autor con Laravel, GitHub y Figma; y (4) que el sistema pueda operar con un modelo de lenguaje local si la empresa lo exige; por rendimiento, el modelo principal está en la nube (Capítulo \ref{chap:implementacion}). La Tabla \ref{tab:dis_herramientas} presenta cada herramienta, las alternativas comparables y la razón de la elección.

{\small
\begin{longtable}{|p{2.4cm}|p{2.9cm}|p{8.3cm}|}
\caption{Herramientas elegidas, alternativas y justificación.}
\label{tab:dis_herramientas}\\
\hline
\textbf{Herramienta} & \textbf{Alternativas} & \textbf{Por qué se eligió} \\
\hline
\endfirsthead
\hline
\textbf{Herramienta} & \textbf{Alternativas} & \textbf{Por qué se eligió} \\
\hline
\endhead
Laravel 13.4 (PHP 8.3) & Django, Node.js/Express, Spring Boot & Incluye ORM, migraciones versionadas, \textit{middleware}, validación, tareas programadas y soporte de pruebas en un solo marco; el autor ya lo conocía. PHP se ejecuta en casi cualquier alojamiento de bajo costo. \\ \hline
React 19 + Inertia.js 2 & Blade, Livewire, React con API REST aparte & React permite construir una interfaz interactiva (arrastrar y soltar, escáner, chat) con un ecosistema amplio de componentes; Inertia evita duplicar la capa de API para las vistas y conserva el enrutamiento y la autenticación de Laravel. \\ \hline
TypeScript 5 & JavaScript & Documenta la forma de los datos que llegan del servidor y permite que el editor señale errores; la verificación estricta en la integración continua queda como trabajo futuro. \\ \hline
Tailwind CSS 4 & Bootstrap, CSS propio & Estilos por utilidades que permiten traducir el prototipo de Figma a componentes sin mantener hojas de estilo separadas por vista, y definir temas mediante variables. \\ \hline
MySQL 8.0 & PostgreSQL, SQLite & Base de datos relacional con transacciones ACID e integridad referencial, necesarias para que un despacho FEFO que toca varios lotes quede completo o no quede. Disponible en el servidor de producción sin costo. En pruebas automatizadas se usa SQLite en memoria. \\ \hline
Ollama y proveedores compatibles con la API de OpenAI & Un único proveedor comercial & Ollama permite ejecutar modelos abiertos en el propio servidor, sin costo por consulta, o en Ollama Cloud, que recibe la pregunta y el contexto filtrado. El diseño admite además proveedores externos configurables desde la aplicación, para elegir según disponibilidad y calidad. \\ \hline
GitHub y GitHub Projects & GitLab, Jira, Trello & Repositorio, historial de cambios, automatización de despliegue y tablero de tareas en una sola plataforma gratuita, que el autor ya usaba. \\ \hline
Figma & Adobe XD, Penpot & Prototipado de alta fidelidad en el navegador, con plan gratuito; el autor ya lo usaba. \\ \hline
Playwright 1.59 y PHPUnit & Cypress, Selenium, Pest & PHPUnit viene integrado con Laravel; Playwright automatiza el navegador real para pruebas de extremo a extremo. \\ \hline
Oracle Cloud (capa Always Free) & AWS, Google Cloud, Azure & Instancia ARM con 4 OCPU y 24 GB de RAM sin costo recurrente \cite{oci_always_free}, suficiente para la aplicación, la base de datos y modelos de lenguaje pequeños. \\ \hline
\end{longtable}
}

\section{Modelo de datos}
\label{sec:dis_datos}

El esquema se define mediante migraciones de Laravel, de modo que cada cambio en la base de datos queda versionado junto con el código. La Tabla \ref{tab:dis_tablas} agrupa las tablas del dominio de Pymetory.

{\small
\begin{longtable}{|p{3.2cm}|p{10.8cm}|}
\caption{Tablas del dominio de Pymetory.}
\label{tab:dis_tablas}\\
\hline
\textbf{Grupo} & \textbf{Tablas y propósito} \\
\hline
\endfirsthead
\hline
\textbf{Grupo} & \textbf{Tablas y propósito} \\
\hline
\endhead
Inventario & \texttt{materials} (catálogo: código, nombre, unidad, categoría, stock mínimo), \texttt{lotes} (número de lote, cantidad, costo unitario, fecha de vencimiento, estado \textit{active}/\textit{quarantined}/\textit{consumed}, material y bodega), \texttt{bodegas} (ubicaciones, capacidad, estado y ruta o enlace de su imagen de fondo), \texttt{transferencias} (traslados entre bodegas). \\ \hline
Kardex & \texttt{movimientos}: lote, usuario, tipo (entrada o salida), cantidad, motivo, descripción y fecha. \\ \hline
Clasificación & \texttt{tags} y \texttt{material\_tag} (etiquetas de materiales), \texttt{custom\_field\_definitions} (campos personalizados). \\ \hline
Compras & \texttt{vendors} (proveedores), \texttt{purchase\_orders} y \texttt{purchase\_order\_items} (órdenes de compra y cantidades recibidas). \\ \hline
Seguridad y auditoría & \texttt{users} (con rol \textit{admin} u \textit{operario}), \texttt{audit\_log} (cambios administrativos), \texttt{api\_keys} (claves de proveedores cifradas). \\ \hline
Asistente & \texttt{chat\_histories} (conversaciones por sesión), \texttt{llm\_providers} (proveedores de modelos configurados), \texttt{llm\_evaluaciones} y \texttt{llm\_precios} (precisión medida y costo de cada modelo), \texttt{settings} (parámetros clave-valor, como el umbral FEFO de 15 días). \\ \hline
Operación & \texttt{notifications} (alertas), \texttt{kanban\_items} (tablero de actividades). \\ \hline
\end{longtable}
}

\subsection{Diseño del Kardex}

La entidad central del diseño es el lote, no el material: la existencia de un material es la suma de las cantidades de sus lotes activos, y cada lote conoce su fecha de vencimiento. Esto es lo que permite aplicar FEFO. Las operaciones que cambian la cantidad de un lote desde el inventario ---ingreso, consumo, despacho por escaneo, transferencia y ajuste de conciliación--- registran un movimiento en la tabla \texttt{movimientos} con el usuario autenticado, el tipo, la cantidad, el motivo y la marca de tiempo. Las operaciones que tocan más de un lote (por ejemplo, un despacho FEFO que agota un lote y continúa con el siguiente, o una transferencia) se ejecutan dentro de una transacción de base de datos, para que el Kardex nunca quede con movimientos a medias.

Por regla de la dirección, el Kardex es de solo inserción, y esa regla se garantiza en dos niveles. En la aplicación, el modelo \texttt{Movimiento} rechaza cualquier intento de modificar o eliminar un registro existente. En la base de datos, las llaves foráneas de \texttt{movimientos} hacia \texttt{lotes} y \texttt{users} impiden borrar un lote o un usuario que tenga historial. Las correcciones se hacen registrando un nuevo movimiento de ajuste con su justificación.

\section{Diseño del módulo de consulta en lenguaje natural}
\label{sec:dis_rag}

El asistente sigue la arquitectura de generación aumentada por recuperación (RAG) \cite{rag_lewis2020}: antes de llamar al modelo de lenguaje, el sistema recupera de la base de datos la información pertinente y la entrega al modelo como contexto, para que la respuesta se base en el inventario real y no en conocimiento general del modelo. Es una variante estructurada: la recuperación son consultas SQL parametrizadas elegidas según la intención, no una búsqueda semántica. El diseño tiene tres etapas:

\begin{enumerate}
    \item \textbf{Clasificación de la intención} (\texttt{intencion}). Un modelo de lenguaje asigna la pregunta a una de trece intenciones o a una general (Tabla \ref{tab:dis_intenciones}); las expresiones regulares de \texttt{classifyQuery} quedan como reserva cuando el modelo no responde o contesta algo fuera de las categorías, y un número de lote se reconoce siempre de forma exacta. Además, las palabras de la pregunta se comparan con los nombres y descripciones de los materiales para filtrar el contexto al material consultado.
    \item \textbf{Construcción del contexto} (\texttt{buildRagContext}). Según la intención, se consulta la base de datos (lotes, movimientos, bodegas, valorización) y se arma un texto con los registros encontrados y el estado de las bodegas. Solo se consideran lotes activos con existencias, el texto empieza con la fecha del día y cada lote indica cuántos días le faltan para vencer, de modo que el modelo pueda responder preguntas como ``¿qué vence esta semana?'' sin calcular fechas por su cuenta.
    \item \textbf{Generación de la respuesta}. El contexto, la instrucción de sistema (editable desde la configuración) y la pregunta se envían al proveedor de modelo de lenguaje configurado; la respuesta se muestra al usuario y se guarda en \texttt{chat\_histories}.
\end{enumerate}

\begin{table}[!htbp]
\centering
\caption{Intenciones reconocidas por el clasificador del asistente.}
\label{tab:dis_intenciones}
\begin{tabular}{|l|p{5.3cm}|p{5.5cm}|}
\hline
\textbf{Intención} & \textbf{Palabras que la activan en los patrones de reserva} & \textbf{Contexto que se recupera} \\
\hline
\texttt{stock\_check} & cuántos hay, stock, cantidad, existencias & Lotes del material, por cantidad \\ \hline
\texttt{critical\_alerts} & crítico, por vencer, alerta, urgente & Lotes dentro del umbral FEFO \\ \hline
\texttt{expiration} & vence, fecha de vencimiento, expira & Lotes por fecha de vencimiento \\ \hline
\texttt{location} & dónde, ubicación, bodega, almacén & Materiales por bodega \\ \hline
\texttt{valuation} & valor, precio, costo & Cantidad por costo unitario \\ \hline
\texttt{movements} & entró, salió, movimiento, historial, kardex & Movimientos con usuario y motivo \\ \hline
\texttt{batch\_info} & lote HAR-01-\ldots & Detalle de un lote \\ \hline
\texttt{summary} & resumen, todo, general, panorama & Resumen del inventario \\ \hline
\texttt{conciliation} & conciliar, ajuste, diferencia, descuadre & Datos para conciliación \\ \hline
\texttt{quarantine} & cuarentena, retenido, no conforme & Lotes en cuarentena \\ \hline
\texttt{low\_stock} & mínimo, hace falta pedir, reponer, agotado & Insumos bajo su stock mínimo \\ \hline
\texttt{forecast} & cuándo se acaba, para cuántos días alcanza, punto de reorden & Proyección: consumo diario, cobertura y punto de reorden \\ \hline
\texttt{consumption} & consumo, cuánto gastamos, por qué bajó & Consumo por periodo y por motivo, comparado con el periodo anterior \\ \hline
\texttt{general} & (ninguna de las anteriores) & Lotes en orden FEFO \\ \hline
\end{tabular}
\end{table}

\subsection{Proveedores de modelo de lenguaje}

El controlador define un mapa de proveedores intercambiables ---Ollama en el propio servidor, OpenCode y cualquier API compatible con la de OpenAI--- y el administrador elige el proveedor y el modelo desde la configuración, con una palanca que muestra la precisión medida y el costo de cada modelo, y con claves de API almacenadas cifradas. Este diseño evita depender de un único proveedor: si el servicio externo no está disponible se puede cambiar a un modelo local, si el proveedor elegido falla responde un modelo rápido en la nube y, después, uno local, y si ningún modelo responde el sistema devuelve en texto plano un listado de lotes en orden FEFO en lugar de un error. La justificación de esta estructura desde el punto de vista de los patrones de diseño se presenta en el capítulo de Implementación.

\section{Diseño de la interfaz de usuario}
\label{sec:dis_ui}

La interfaz se prototipó en Figma. El commit del 14 de abril de 2026 (\texttt{38377c4}, ``Reconstrucción Figma de Alta Fidelidad (18 vistas)'') trasladó a React las 18 vistas de alta fidelidad de ese prototipo; por eso varios componentes conservan el prefijo \texttt{Figma} en su nombre (\texttt{FigmaTablero}, \texttt{FigmaLabelPrint}, \texttt{FigmaQRScanner}, entre otros).

Los principios que guiaron el diseño de la interfaz fueron:

\begin{itemize}
    \item \textbf{Navegación por tareas.} La barra lateral agrupa los módulos en tres secciones ---Principal (Tablero, Inventario, Buscar), Análisis (Asistente, Reportes, Reabastecimiento, Kardex, Etiquetas, Escáner) y Gestión (Historial QR, Transferencias, Órdenes de compra, Conteo físico, Imprimir etiquetas, Alertas)--- y el tablero ofrece accesos rápidos a las acciones más frecuentes (nuevo insumo, escanear, reporte, consumo rápido).
    \item \textbf{Lo urgente primero.} El tablero muestra en la parte superior los indicadores que requieren acción (lotes críticos por vencimiento, existencias, valorización) y, debajo, la actividad reciente del Kardex y la ocupación de las bodegas.
    \item \textbf{Uso en bodega desde el celular.} El diseño es adaptable a pantallas pequeñas: el menú lateral se reemplaza por una barra inferior con las cuatro acciones más usadas (tablero, inventario, escanear y asistente) y un panel con el resto de módulos, y el escáner usa la cámara del dispositivo, para que el personal pueda registrar movimientos sin ir a un computador.
    \item \textbf{Temas visuales.} Los colores se definen como variables de CSS, lo que permitió ofrecer seis temas seleccionables (\textit{midnight-luxe}, \textit{carbon-amber}, \textit{nordic-steel}, \textit{obsidian-teal}, \textit{royal-plum} y \textit{paper-light}), incluido uno claro. Cada tema define un pequeño conjunto de colores base (fondo, panel, borde, texto y acento) y a partir de ellos redefine la paleta completa que usan los componentes; así un tema no puede alterar a otro. Las imágenes, como las de las bodegas, provienen de la base de datos (archivo subido o enlace), no de rutas fijas en el código.
\end{itemize}
